% =============================================================================
% proofs/07_detection.tex
% Detection, silence, and bidder entry at fixed policies.
% Fragment, to be \input by appendix.tex.
% =============================================================================

\section{Detection: the entry identity, the tape, who gets caught and silence}
\label{sec:detection}

This section connects the engagement premium to bidder entry and separates the two ways public
information identifies a Voice type. The filing acts through the disclosure cell. The tape acts
through flow values that reveal a building purchase.
At any date $d$, write $D^d$ for the indicator that the filing has landed by $d$, and write
$\mathcal C_F^d=\{D^d=1\}$ and $\mathcal C_P^d=\{D^d=0\}$.

\csname unlabelledtrue\endcsname
\begin{lemma}[Entry identity]
\label{lem:entry-identity}
Let $\mathcal I_H$ be the public information set at the control node, let
$\pi(\mathcal I_H)=\Prb(a=1\mid\mathcal I_H)$, and let
$p(\mathcal I_H)$ be the bidder-entry probability. The premium kernel is
$h(\mathcal I_H)=\pi(\mathcal I_H)p(\mathcal I_H)$. Assume:
\begin{enumerate}[label=(E-\arabic*),leftmargin=3.2em,itemsep=3pt]
\item \emph{Public measurability.} The price and the entry probability are functions of the public
  information set, and the engagement indicator is a function of the type. Thus
  $P=P(\mathcal I_H)$ and $p=p(\mathcal I_H)$ are $\mathcal I_H$-measurable, while
  $a=a(\theta)$. The first assertion is the pinned-kernel content of (S8); the last follows from
  the finite menu and cutoff selection map in (S2) and (S3).
\item \emph{The Voice event and the public cells.} The menu assigns $a=1$ exactly to Voice, and
  $a=0$ to Exit and Hold. The cell indicator $D$ is public by (S5) and (S6), so
  $\mathcal C_F=\{D=1\}$ and $\mathcal C_P=\{D=0\}$ belong to $\mathcal I_H$.
\item \emph{Finite premium and fixed policies.} The premium satisfies (S9), liquidity enters only
  through the noise law as in (S10), and the menu, cutoff vector, and execution policies are fixed
  as in (S11).
\end{enumerate}
Write $\omega_V=\Prb(a=1)$ for the Voice mass. If $\omega_V>0$, define
$e_V=\Eop[p(\mathcal I_H)\mid a=1]$. Then
\begin{equation}
\label{eq:entry-identity}
  \Dact=\Delta_m\,\omega_V e_V .
\end{equation}
For either cell $\mathcal C\in\{\mathcal C_F,\mathcal C_P\}$ with
$\Prb(\mathcal C)>0$ and $\Prb(a=1,\mathcal C)>0$, put
\[
  \omega_{V\mid\mathcal C}=\Prb(a=1\mid\mathcal C),
  \qquad
  e_{V\mid\mathcal C}=\Eop[p(\mathcal I_H)\mid a=1,\mathcal C].
\]
Its cell average obeys the same identity,
\begin{equation}
\label{eq:entry-identity-cell}
  M_{\mathcal C}
  :=\Delta_m\Eop[\pi(\mathcal I_H)p(\mathcal I_H)\mid\mathcal C]
  =\Delta_m\,\omega_{V\mid\mathcal C}e_{V\mid\mathcal C}.
\end{equation}
If the Voice mass is zero, the corresponding premium is zero and the conditional entry
probability is left undefined. At fixed policies, $\omega_V$ is independent of both $\kappa$ and
the disclosure rule.
\end{lemma}
\csname unlabelledfalse\endcsname

\begin{proof}
By (E-1), $p(\mathcal I_H)$ is $\mathcal I_H$-measurable, and (S8) makes it bounded. The defining
property of conditional expectation therefore gives
\[
  \Eop[\pi(\mathcal I_H)p(\mathcal I_H)]
  =\Eop\!\left[\Eop[a\mid\mathcal I_H]p(\mathcal I_H)\right]
  =\Eop[a\,p(\mathcal I_H)].
\]
When $\omega_V>0$, conditioning the last expectation on the event $\{a=1\}$ gives
$\Eop[a p]=\omega_V\Eop[p\mid a=1]=\omega_Ve_V$. Multiplication by $\Delta_m$ and (S9) yield
\eqref{eq:entry-identity}. If $\omega_V=0$, then $a=0$ almost surely, so both
$\Eop[a p]$ and the engagement premium vanish.

Now fix a cell $\mathcal C$ of positive mass. Its indicator is $\mathcal I_H$-measurable by
(E-2). Applying the same calculation after multiplication by $\ind_{\mathcal C}$ gives
\[
  \Eop[\pi p\,\ind_{\mathcal C}]
  =\Eop[a p\,\ind_{\mathcal C}].
\]
Divide by $\Prb(\mathcal C)$. If $\Prb(a=1,\mathcal C)>0$, conditioning the right side first on
$\{a=1\}\cap\mathcal C$ gives \eqref{eq:entry-identity-cell}. If that joint mass is zero, the
right side is zero, hence so is the cell premium.

It remains to check the invariance of the unconditional Voice mass. By (S3) and (S7),
$a=a_{j(s)}$. Condition (S11) holds the menu and the cutoff map $j(\cdot)$ fixed across rules and
in $\kappa$, so neither the threshold margin nor the window margin enters this indicator.
Condition (S10) puts $\kappa$ only in the law of the noise marks, while the law of the signal is
unchanged by (S1). The law of $a_{j(s)}$, and hence $\omega_V$, is therefore common to every rule
and every $\kappa$ in a fixed-policy comparison.
\end{proof}

\csname unlabelledtrue\endcsname
\begin{lemma}[Detection]
\label{lem:detection}
Fix a date $d\in\{0,\ldots,H\}$ and a pooled history of positive probability. Assume:
\begin{enumerate}[label=(D-\arabic*),leftmargin=3.2em,itemsep=3pt]
\item \emph{Voice marks.} The order mark is $g_e(\theta)\in\{0,1\}$ and strategic order flow is
  $q_e(\theta)=2\bar z\,g_e(\theta)$. Only Voice plans carry positive marks. Exit and Hold have
  $g_e=0$ at every date, and a Voice building round has $g_e=1$.
\item \emph{Common noise channel.} Condition (S12) holds. In particular,
  $X_e=q_e+z_e$, the noise marks are independent across dates and independent of the type, and
  $z_e$ equals $-\bar z$, $0$, or $+\bar z$ with probabilities
  $\kappa/2$, $1-\kappa$, and $\kappa/2$, where $\kappa\in[0,1]$.
\item \emph{Public observation.} Conditions (S6), (S7), (S8), and (S13) hold, so the pooled flow
  history is public, its posterior is the pinned conditional probability, and the date-$d$ pooled
  cell is a type event.
\end{enumerate}
The flow values $2\bar z$ and $3\bar z$ reveal a building mark. On any positive-probability
pooled history containing either value, the engagement posterior is one. The values $-\bar z$
and $0$ reveal a zero mark whenever those values have positive probability. The value $+\bar z$
is the only ambiguous value: in the round-level channel support at $\kappa>0$, it can arise from a
building purchase with a noise sale or from a zero mark with a noise purchase. At $\kappa=0$ the
ambiguous value is not observed.

For a Voice type, let
\[
  N_d(\theta)=\sum_{e=0}^{d}g_e(\theta)
\]
be the number of building rounds through that date. The probability that no building mark has
been revealed on the tape is
\begin{equation}
\label{eq:detection-silence}
  \Prb(\text{no building mark revealed through }d\mid\theta)
  =\left(\frac{\kappa}{2}\right)^{N_d(\theta)}.
\end{equation}
The right side is the product over building rounds, with an empty product equal to one.
If the type remains pooled at date $d$, this event is a silent history and the same formula is its
silent probability.
\end{lemma}
\csname unlabelledfalse\endcsname

\begin{proof}
If $g_e=0$, then $q_e=0$, so the three possible flow values are
$-\bar z$, $0$, and $+\bar z$. If $g_e=1$, then $q_e=2\bar z$, so the three possible values are
$+\bar z$, $2\bar z$, and $3\bar z$. The two supports meet only at $+\bar z$. This proves the
stated classification.

The event $\{X_e\in\{2\bar z,3\bar z\}\}$ implies $g_e=1$. By (D-1), it therefore implies that
the selected plan is Voice and $a=1$. On a pooled history of positive probability that contains
such a value, Bayes' rule assigns probability one to $a=1$. Thus
$\pi(\mathcal I_d)=1$ on that history.

Fix a Voice type. On each of its $N_d(\theta)$ building rounds, the only flow value
that does not reveal the positive mark is $+\bar z$. That value occurs on a building round
exactly when $z_e=-\bar z$, an event of probability $\kappa/2$. Zero-mark rounds cannot produce
$2\bar z$ or $3\bar z$ and hence cannot reveal a building mark. Independence across dates in
(S12) now gives the product in \eqref{eq:detection-silence}. The event $D^d=0$ is a function of
the type under (S7), so restricting the formula to pooled types does not change the noise product
law.
\end{proof}

\csname unlabelledtrue\endcsname
\begin{lemma}[Upper set]
\label{lem:upper-set}
Fix a threshold margin $\tau$, a window margin $T$, and a date
$d\in\{0,\ldots,H\}$. Assume:
\begin{enumerate}[label=(U-\arabic*),leftmargin=3.2em,itemsep=3pt]
\item \emph{Monotone target.} The Voice target $b^*(s)\ge b_0$ is weakly increasing in the
  signal.
\item \emph{Monotone accumulation length.} The accumulation length
  $n(s)\in\{1,\ldots,H+1\}$ is weakly decreasing in the signal.
\item \emph{Ordered menu form.} Voice is selected on an upper interval of the signal line, and
  its stake and mark paths are
  \[
    B_V(s,e)=b_0+\bigl(b^*(s)-b_0\bigr)
      \min\!\left\{1,\frac{e+1}{n(s)}\right\},
    \qquad
    g_e(s)=\ind\{e<n(s)\}.
  \]
\item \emph{Timing and tape.} Conditions (S4), (S5), and (S7) hold, together with (D-1) and the
  common noise channel (S12).
\end{enumerate}
Then the date-$d$ flagged signal set
\begin{equation}
\label{eq:upper-flagged-set}
  \mathcal F_d
  =\{s:j(s)=\mathrm{Voice},\ c(s;\tau)+T\le d\}
\end{equation}
is an upper set. If $s_2\ge s_1$ are Voice signals, then
\[
  c(s_2;\tau)\le c(s_1;\tau),
  \qquad
  N_d(s_2)\le N_d(s_1),
  \qquad
  N_d(s)=\min\{n(s),d+1\}.
\]
Consequently, for prior-almost every pair of Voice types $\theta_F$ and $\theta_P$ with
$s(\theta_F)\in\mathcal F_d$ and
$s(\theta_P)\notin\mathcal F_d$,
\begin{equation}
\label{eq:upper-silence-order}
  \Prb(\text{no building mark revealed through }d\mid\theta_F)
  \ge
  \Prb(\text{no building mark revealed through }d\mid\theta_P).
\end{equation}
The inequality is strict when $\kappa>0$ and $N_d(\theta_F)<N_d(\theta_P)$. Thus the rule flags from
the Voice types with weakly fewer building rounds, which occupy the weakly highest-probability end
among Voice types for escaping tape identification.
\end{lemma}
\csname unlabelledfalse\endcsname

\begin{proof}
Take Voice signals $s_2\ge s_1$. By (U-1),
$b^*(s_2)-b_0\ge b^*(s_1)-b_0\ge0$. By (U-2),
$n(s_2)\le n(s_1)$, and therefore, at every date $e$,
\[
  \min\!\left\{1,\frac{e+1}{n(s_2)}\right\}
  \ge
  \min\!\left\{1,\frac{e+1}{n(s_1)}\right\}.
\]
Both factors in the product defining the increment above $b_0$ are nonnegative and weakly
larger at $s_2$. The path formula in (U-3) gives
$B_V(s_2,e)\ge B_V(s_1,e)$ for every $e$.

Under (S5), $c(s;\tau)$ is the first date at which this path reaches $\tau$. If the lower signal
has crossed by a date, the higher signal has crossed by that date as well. Hence
$c(s_2;\tau)\le c(s_1;\tau)$. If $s_1\in\mathcal F_d$, then (U-3) keeps every higher signal in
Voice and the crossing inequality gives
$c(s_2;\tau)+T\le c(s_1;\tau)+T\le d$. Thus $s_2\in\mathcal F_d$, proving that
$\mathcal F_d$ is an upper set. The path is fixed before flow is observed by (S7), so this
ordering is not altered by realised noise or prices.

The mark formula in (U-3) makes the mark equal to one at each date
$e<\min\{n(s),d+1\}$ and none after it, so
$N_d(s)=\min\{n(s),d+1\}$. This count is weakly decreasing in $s$ by (U-2). Since
$\mathcal F_d$ is upper, a flagged signal cannot lie below an unflagged Voice signal. Thus, for the
types in the statement, $N_d(\theta_F)\le N_d(\theta_P)$. Lemma~\ref{lem:detection} gives the two
no-revelation probabilities as $(\kappa/2)^{N_d(\theta_F)}$ and
$(\kappa/2)^{N_d(\theta_P)}$. The base lies in $[0,1/2]$, so the smaller
exponent gives the weakly larger probability. It gives a strict inequality when the base is
positive and the exponents differ.
\end{proof}

\csname unlabelledtrue\endcsname
\begin{lemma}[Silence]
\label{lem:silence}
Fix a date $d\in\{0,\ldots,H\}$ and compare rules at common fixed policies and common $\kappa$.
Either $r_+=(\tau',T)$ and $r_-=(\tau,T)$ with $b_0<\tau'<\tau$, or
$r_+=(\tau,T')$ and $r_-=(\tau,T)$ with $T'<T$. Let
\[
  A=\mathcal C_{P,r_+}^{d},
  \qquad
  R=\mathcal C_{P,r_-}^{d}\setminus\mathcal C_{P,r_+}^{d}.
\]
Assume:
\begin{enumerate}[label=(L-\arabic*),leftmargin=3.2em,itemsep=3pt]
\item \emph{Nested upper removal.} Conditions (S4), (S5), (S7), and (S11), together with
  (U-1) to (U-3), imply
  $\mathcal C_{P,r_+}^{d}\subseteq\mathcal C_{P,r_-}^{d}$. The set $R$ has positive prior mass,
  and $A$ and $R$ are $\sigma(s)$-measurable at fixed policies by Step~4 of
  Lemma~\ref{lem:partition}. The set $R$ lies in $\{a=1\}$ and is upper relative to $A$: for
  prior-almost every
  $\theta_R\in R$ and $\theta_A\in A$, $s(\theta_R)\ge s(\theta_A)$.
\item \emph{Common pooled likelihood.} Conditions (S12) and (S13) hold. On $A$, the two rules
  have the same type-conditional law of the displayed pooled flow history.
\item \emph{Monotone signal information.} The conditional mean
  $\mu_v(s)=\Eop[v\mid s]$ exists and is weakly increasing in $s$, and
  $v$ is conditionally independent of pooled flow given $s$. In the Gaussian signal model,
  $\mu_v(s)=\mu_v+\beta(s-\mu_v)$ with $\beta>0$, so this ordering is strict.
\item \emph{Common positive-probability history.} The displayed pooled flow history
  $x=(x_0,\ldots,x_d)$, with all public filing coordinates equal to zero, has positive
  probability under both rules.
\end{enumerate}
Let
\[
  \pi_r(x)=\Prb(a=1\mid X_{0:d}=x,D_r^d=0),
  \qquad
  \widehat v_r(x)=\Eop[v\mid X_{0:d}=x,D_r^d=0].
\]
Then
\begin{equation}
\label{eq:silence-posterior-order}
  \pi_{r_+}(x)\le\pi_{r_-}(x),
  \qquad
  \widehat v_{r_+}(x)\le\widehat v_{r_-}(x).
\end{equation}
The first inequality is strict exactly when the looser posterior assigns positive probability to
$R$ and $\pi_{r_+}(x)<1$. The second is strict when the looser posterior assigns positive
probability to $R$ and the posterior mean of $\mu_v(s)$ on $R$ is strictly larger than its
posterior mean on $A$. In particular, strict monotone signal information and strictly separated
posterior signal supports give strictness when the looser posterior assigns positive probability
to $R$. The conclusions apply to every silent history that satisfies (L-4), as well as to every
other common pooled history of positive probability.
\end{lemma}
\csname unlabelledfalse\endcsname

\begin{proof}
Nestedness partitions the looser pooled type set into $A$ and $R$. Condition (S11) gives both rules
the same prior over types. By (L-2), the likelihood of
the common flow vector $x$ given a type is the same under both rules. Conditioning the looser
posterior further on $A$ therefore gives the tighter posterior. Put
\[
  \lambda_x=\Prb(R\mid X_{0:d}=x,D_{r_-}^d=0).
\]
Assumption (L-4) gives positive posterior mass to $A$, so $0\le\lambda_x<1$. If
$\lambda_x=0$, the two posteriors coincide and both inequalities in
\eqref{eq:silence-posterior-order} are equalities. Suppose now that $\lambda_x>0$.

Every type in $R$ is engaged by (L-1). The looser engagement posterior is therefore the mixture
\[
  \pi_{r_-}(x)
  =(1-\lambda_x)\pi_{r_+}(x)+\lambda_x.
\]
Since $\pi_{r_+}(x)\le1$, this proves the first inequality. The displayed mixture also shows that
it is strict exactly under the conditions stated.

Let
\[
  \widehat v_R(x)=\Eop[v\mid X_{0:d}=x,\theta\in R].
\]
The same posterior decomposition gives
\begin{equation}
\label{eq:silence-v-mixture}
  \widehat v_{r_-}(x)
  =(1-\lambda_x)\widehat v_{r_+}(x)+\lambda_x\widehat v_R(x).
\end{equation}
By (L-3), conditioning on the flow adds no information about $v$ after conditioning on $s$.
Thus the two component means in \eqref{eq:silence-v-mixture} are averages of $\mu_v(s)$ under
their respective posterior signal laws. Assumption (L-1) puts every signal in $R$ weakly above
every signal in $A$, and (L-3) makes $\mu_v$ weakly increasing. Hence
$\widehat v_R(x)\ge\widehat v_{r_+}(x)$. Substitution into
\eqref{eq:silence-v-mixture} proves the second inequality. It is strict exactly when
$\lambda_x>0$ and the component means differ, with the stated signal conditions sufficient for
that difference.
\end{proof}

\csname unlabelledtrue\endcsname
\begin{lemma}[Entry against the undetected]
\label{lem:entry-undetected}
Let $\Pi(\widehat v,\pi)$ be the unique pricing root of
Lemma~\ref{lem:pricing-root}, and define the entry probability at that root by
$e(\widehat v,\pi)=p(\Pi(\widehat v,\pi),\pi)$. Fix an interval
$\mathcal V\subset\mathbb R$ of fundamental posterior means. Assume:
\begin{enumerate}[label=(I-\arabic*),leftmargin=3.2em,itemsep=3pt]
\item \emph{Unique pinned root.} The hypotheses of Lemma~\ref{lem:pricing-root} hold, including
  $m_0\ge0$, $\Delta_m>0$, $\Delta_V\ge0$, $\sigma_\xi>0$, and uniqueness of the finite fixed
  point. The pinned pricing rule is the one in (S8).
\item \emph{Entry derivatives.} For every finite $P$ and $\pi\in[0,1]$, entry has the normal form
  \[
    p(P,\pi)=1-\Phi\!\left(
      \frac{P+K+m_0+\pi\Delta_m-\bar S}{\sigma_\xi}\right),
  \]
  so it is continuously differentiable with
  $\partial_Pp<0$ and $\partial_\pi p<0$.
\item \emph{Added engagement-price condition.} At every root
  $P=\Pi(\widehat v,\pi)$ with $(\widehat v,\pi)\in\mathcal V\times[0,1]$, write
  $m(\pi)=m_0+\pi\Delta_m$ and $o(P,\pi)=p(P,\pi)/(1-p(P,\pi))$. Assume
  \begin{equation}
  \label{eq:entry-price-condition}
    \Delta_V+\Delta_m o(P,\pi)+m(\pi)\,\partial_\pi o(P,\pi)>0.
  \end{equation}
  Equivalently, with
  $t=(P+K+m(\pi)-\bar S)/\sigma_\xi$ and standard normal density $\phi$,
  \[
    \Delta_V+\Delta_m o(P,\pi)
    >
    \frac{m(\pi)\Delta_m\phi(t)}{\sigma_\xi\Phi(t)^2}.
  \]
\item \emph{Added history-comparison condition.} For the history comparison, the silent pooled
  history has engagement posterior $\pi_U\le1$. A tape-detected pooled history with the same fundamental
  posterior mean has engagement posterior one. The comparison is nondegenerate when $\pi_U<1$.
\end{enumerate}
Then $\Pi(\widehat v,\pi)$ is strictly increasing in both $\widehat v$ and $\pi$ on the stated
domain, while $e(\widehat v,\pi)$ is strictly decreasing in both. Hence, holding the fundamental
posterior mean fixed,
\begin{equation}
\label{eq:entry-undetected-gap}
  e(\widehat v,\pi_U)\ge e(\widehat v,1),
  \qquad \widehat v\in\mathcal V.
\end{equation}
The inequality is strict when $\pi_U<1$ and is an equality when $\pi_U=1$.
\end{lemma}
\csname unlabelledfalse\endcsname

\begin{proof}
The pricing fixed point can be written as the root of
\begin{equation}
\label{eq:entry-root-residual}
  \mathcal R(P;\widehat v,\pi)
  =P-\widehat v-\pi\Delta_V-m(\pi)o(P,\pi)=0.
\end{equation}
By (I-2), $p\in(0,1)$, $\partial_Pp<0$, and the odds map is strictly increasing in $p$.
Therefore $\partial_Po<0$. Since $m(\pi)\ge0$,
\begin{equation}
\label{eq:entry-root-slope}
  \partial_P\mathcal R=1-m(\pi)\partial_Po>0.
\end{equation}
Lemma~\ref{lem:pricing-root} supplies the unique root. Equation
\eqref{eq:entry-root-slope} lets the implicit function theorem differentiate that root.
Because $\partial_{\widehat v}\mathcal R=-1$,
\[
  \partial_{\widehat v}\Pi
  =-\frac{\partial_{\widehat v}\mathcal R}{\partial_P\mathcal R}
  =\frac{1}{\partial_P\mathcal R}>0.
\]
For the engagement posterior,
\[
  \partial_\pi \mathcal R
  =-\Delta_V-\Delta_mo-m(\pi)\partial_\pi o.
\]
Condition \eqref{eq:entry-price-condition} says exactly that this derivative is negative. Hence
\[
  \partial_\pi\Pi
  =-\frac{\partial_\pi \mathcal R}{\partial_P\mathcal R}>0.
\]
Continuity of the unique root, supplied by Lemma~\ref{lem:pricing-root}, extends these monotone
comparisons to the endpoints of the posterior interval.

Apply the chain rule to entry at the root:
\[
  \partial_{\widehat v}e
  =(\partial_Pp)(\partial_{\widehat v}\Pi)<0,
  \qquad
  \partial_\pi e
  =(\partial_Pp)(\partial_\pi\Pi)+\partial_\pi p<0.
\]
Both signs use all the derivative restrictions in (I-2); the second also uses (I-3). Under (I-4),
decrease in $\pi$ at fixed $\widehat v$ gives \eqref{eq:entry-undetected-gap}, strictly when
$\pi_U<1$. When $\pi_U=1$, both arguments of $e$ are the same.
\end{proof}
